\chapter{ԵԶՐԱԿԱՑՈՒԹՅՈՒՆՆԵՐ ԵՎ ԱՌԱՋԱՐԿՈՒԹՅՈՒՆՆԵՐ}
\sloppy

Սույն ավարտական աշխատանքում ուսումնասիրվեց և գործնականորեն իրականացվեց «Թեմատիկ վեբ կայքի ստեղծում առաջատար վեբ ծրագրավորմամբ» թեման՝ անձնական նպատակների հետևման առարկայական ուղղությամբ։ Աշխատանքի սկզբում ձևակերպված հիմնական խնդիրն այն էր, որ օգտատերերը հաճախ ստիպված են նպատակների սահմանումը, ենթանպատակների պլանավորումը, առաջընթացի գրանցումը, հիշեցումները և արդյունքների վերլուծությունը կատարել տարբեր գործիքներում։ Այդ մասնատվածությունը մեծացնում է ճանաչողական ծանրաբեռնվածությունը և նվազեցնում է երկարաժամկետ նպատակների նկատմամբ հետևողականությունը։ Կատարված ուսումնասիրությունը և մշակված վեբ կայքը ցույց տվեցին, որ առաջատար full-stack վեբ ծրագրավորման մոտեցումները կարող են այդ գործառույթները միավորել մեկ հասանելի, պաշտպանված և ընդլայնելի համակարգում։

Առաջին խնդիրը գիտական մոտեցումների գնահատումն էր։ Գրականության վերլուծության արդյունքում պարզվեց, որ անձնական նպատակների արդյունավետ կառավարումը պետք է հիմնվի չափելիության, հստակության, փուլային բաժանման և պարբերական հետադարձ կապի վրա։ SMART մոտեցումը, Goal Setting Theory-ն և OKR տրամաբանությունը ցույց են տալիս, որ նպատակը պետք է լինի ոչ միայն ձևակերպված, այլև չափելի և վերահսկելի։ Ներգրավվածության վերաբերյալ աղբյուրները հաստատեցին, որ self-monitoring-ը, progress feedback-ը, streak-ները և չափավոր հիշեցումները կարող են աջակցել օգտատիրոջ հետևողականությանը, եթե դրանք չեն վերածվում ավելորդ ճնշման։ Այս եզրակացությունը դարձավ ծրագրային պահանջների հիմք՝ նպատակ, ենթանպատակ, առաջընթացի պատմություն, dashboard և analytics գործառույթների համար։

Երկրորդ խնդիրը գործող թվային լուծումների համեմատությունն էր։ Todoist, Habitica, Strides և Goals on Track գործիքների ուսումնասիրությունը ցույց տվեց, որ յուրաքանչյուր լուծում ունի ուժեղ կողմ, սակայն դրանք հաճախ կենտրոնացած են մեկ ուղղության վրա՝ task management, gamification, habit tracking կամ SMART պլանավորում։ Համեմատությունը ապացուցեց, որ բակալավրի աշխատանքի ընտրված թեմայի համար անհրաժեշտ էր ոչ թե կրկնել առկա գործիքներից մեկը, այլ ստեղծել հավասարակշռված թեմատիկ կայք, որտեղ պլանավորումը, առաջընթացի գրանցումը և վերլուծությունը դիտարկվում են որպես նույն գործընթացի հաջորդական փուլեր։

Երրորդ խնդիրը կայքի կառուցվածքի, տվյալների մոդելի և հիմնական գործառույթների մշակումն էր։ Այս խնդիրը իրականացվեց Next.js 15 App Router, TypeScript, Prisma, PostgreSQL, Auth.js v5, Tailwind CSS, shadcn/ui, Recharts և Docker տեխնոլոգիաներով։ Մշակված տվյալների մոդելը ներառում է User, Category, Goal, Milestone և ProgressLog սուբյեկտներ, որոնք միասին ապահովում են անհատականացված նպատակների պահպանում, դասակարգում, ենթանպատակների կառավարում և առաջընթացի պատմություն։ Հավելվածում իրականացվեցին նույնականացում, պաշտպանված էջեր, նպատակների CRUD գործողություններ, ֆիլտրեր, որոնում, dashboard, analytics, export, հիշեցումներ, light/dark թեմա և responsive միջերես։

Չորրորդ խնդիրը օգտագործման հարմարավետության, ֆունկցիոնալ ամբողջականության և օգտատիրոջ հիմնական պահանջներին համապատասխանության գնահատումն էր։ Արդյունքները ցույց տվեցին, որ մշակված կայքը բավարարում է հիմնական պահանջները. օգտատերը կարող է մուտք գործել համակարգ, ստեղծել չափելի նպատակ, բաժանել այն ենթանպատակների, թարմացնել ընթացիկ արժեքը, տեսնել պատմական գրառումները, դիտել ընդհանուր վիճակագրությունը և վերլուծել արդյունքները ըստ կատեգորիաների ու ժամանակային միջակայքերի։ Responsive դասավորությունը և mobile dashboard-ի ստուգումը ցույց տվեցին, որ համակարգը կարող է օգտագործվել ոչ միայն desktop, այլև փոքր էկրանների վրա։

Հինգերորդ խնդիրը առաջարկվող լուծման միջոցով նպատակների կառավարման մասնատվածության նվազեցման ապացուցումն էր։ Մշակված համակարգը այդ խնդիրը լուծում է ինտեգրման սկզբունքով. նույն միջավայրում միավորվում են նպատակների սահմանումը, գործողությունների փուլային բաժանումը, առաջընթացի գրանցումը, հիշեցումները և վերլուծական պատկերները։ Այս կապը պահպանում է աշխատանքում պահանջվող խնդրից դեպի մեթոդ, արդյունք, եզրակացություն և առաջարկություն շղթան. մասնատվածության խնդիրը ուսումնասիրվեց գրականության և համեմատական վերլուծության միջոցով, մեթոդական հիմքը դարձավ իտերատիվ full-stack մշակում, իսկ արդյունքը՝ ամբողջական թեմատիկ վեբ կայք։

Աշխատանքի հիման վրա ձևակերպվում են հետևյալ առաջարկությունները։ Նախ՝ համակարգի հաջորդ տարբերակում նպատակահարմար է ավելացնել անհատականացված շաբաթական վերանայման մեխանիզմ, որը օգտատիրոջը ոչ թե պարզապես կծանուցի վերջնաժամկետի մասին, այլ կցուցադրի շաբաթվա առաջընթացը, ուշացած նպատակները և հաջորդ ամենակարևոր քայլը։ Այս առաջարկությունը հիմնվում է գրականության այն եզրակացության վրա, որ self-monitoring-ը և հետադարձ կապը բարձրացնում են ներգրավվածությունը, ինչպես նաև մշակված dashboard և analytics էջերի արդյունքների վրա։

Երկրորդ՝ անհրաժեշտ է ընդլայնել usability testing-ը իրական օգտատերերի մասնակցությամբ՝ առնվազն ուսանողների և երիտասարդ մասնագետների փոքր խմբով։ Թեստավորումը պետք է չափի, թե որքան արագ են օգտատերերը ստեղծում նպատակ, ավելացնում ենթանպատակ, գրանցում առաջընթաց և գտնում analytics տվյալները։ Այս առաջարկությունը կապված է աշխատանքի չորրորդ խնդրի հետ, քանի որ ներկայիս գնահատումը հիմնված է ֆունկցիոնալ և փորձնական ստուգման վրա, իսկ իրական օգտագործման տվյալները թույլ կտան ավելի ճշգրիտ գնահատել միջերեսի պարզությունը և ճանաչողական ծանրաբեռնվածությունը։

Երրորդ՝ հետագա զարգացման փուլում նպատակահարմար է ավելացնել կարգավորվող push notification-ներ և calendar integration, սակայն դրանք պետք է լինեն օգտատիրոջ կողմից վերահսկվող։ Գրականության վերլուծությունը ցույց տվեց, որ թվային բարեկեցության տեսանկյունից ծանուցումները պետք է աջակցեն ինքնակարգավորմանը, ոչ թե ավելացնեն ուշադրության ծանրաբեռնվածությունը։ Հետևաբար հիշեցումների ընդլայնումը պետք է իրականացվի նախընտրելի հաճախականության, անջատման և նպատակին կապվող իմաստալից հաղորդագրությունների սկզբունքով։

Աշխատանքն ունի նաև սահմանափակումներ։ Մշակված համակարգը դեռ չունի native mobile հավելված, թիմային նպատակների կառավարում, լայնածավալ օգտագործման վիճակագրություն և արտադրական միջավայրում երկարաժամկետ փորձարկում։ Բացի դրանից, գնահատումը հիմնականում կատարվել է ֆունկցիոնալ ամբողջականության, կառուցվածքային համեմատության և responsive ստուգման հիման վրա, այլ ոչ թե մեծ օգտատերային նմուշով։ Այս սահմանափակումները չեն փոխում աշխատանքի հիմնական եզրակացությունը, սակայն դրանք կարևոր են ապագա հետազոտության և զարգացման ուղղությունները սահմանելու համար։

Ապագա զարգացման հիմնական ուղղություններն են՝ native կամ progressive web app տարբերակ, անհատականացված առաջարկությունների համակարգ, թիմային կամ mentor-based նպատակների աջակցություն, իրական օգտագործման analytics, տվյալների ներմուծում և արտահանում տարբեր ձևաչափերով, ինչպես նաև accessibility audit՝ WCAG 2.2 չափանիշներով։ Այս ուղղությունները կարող են ընդլայնել մշակված թեմատիկ կայքը ուսումնական նախագծից դեպի ավելի ամբողջական արտադրողականության հարթակ։

Այսպիսով, աշխատանքը հասավ իր նպատակին. ուսումնասիրվեցին անձնական նպատակների հետևման և առաջատար վեբ ծրագրավորման տեսական ու գործնական հիմքերը, համեմատականորեն գնահատվեցին առկա լուծումները, մշակվեց full-stack թեմատիկ վեբ կայք, և ցույց տրվեց, որ մեկ համակարգում պլանավորման, առաջընթացի հետևման և վերլուծության միավորումը կարող է նվազեցնել նպատակների կառավարման մասնատվածությունը։

\fussy
